\documentclass[12pt,a4paper]{article}

\usepackage[utf8]{inputenc}
\usepackage[T1]{fontenc}
\usepackage[brazil]{babel}
\usepackage{lmodern}
\usepackage{microtype}
\usepackage[a4paper,margin=2.5cm]{geometry}
\usepackage{booktabs}
\usepackage{longtable}
\usepackage{tabularx}
\usepackage{array}
\usepackage{enumitem}
\usepackage[hidelinks]{hyperref}
\usepackage{xcolor}
\usepackage{csquotes}

\newcolumntype{Y}{>{\raggedright\arraybackslash}X}
\setlist[itemize]{noitemsep,topsep=4pt}
\setlist[enumerate]{noitemsep,topsep=4pt}
\hypersetup{pdftitle={Resumo de pesquisa: Vertentes contemporâneas em Educação Ambiental},pdfauthor={}}

\title{Vertentes contemporâneas em Educação Ambiental\\\large Resumo de pesquisa para leitura e estudo}
\author{}
\date{Setembro de 2026}

\begin{document}
\maketitle

\begin{abstract}
Este texto organiza, em forma de resumo de estudo, os achados dos artefatos de pesquisa produzidos sobre vertentes contemporâneas em Educação Ambiental (EA). Ele prioriza seis obras indicadas: Layrargues e Lima; Pelicioni e Philippi Jr.; Santos e Toschi; Carvalho; Guimarães; e Souza. As afirmações atribuídas a autores são paráfrases das fontes ou dos fichamentos gerados. Quando uma relação é construída pela síntese da pesquisa, ela é indicada como \emph{interpretação}. Não se trata de uma apresentação nem de uma nova pesquisa bibliográfica.
\end{abstract}

\tableofcontents
\newpage

\section{Como ler este resumo}

O achado mais importante da pesquisa é que Educação Ambiental não designa uma doutrina única. Ela é um campo plural de práticas, conceitos e disputas. Assim, duas atividades que parecem semelhantes --- por exemplo, uma campanha de reciclagem --- podem se apoiar em diagnósticos muito diferentes da crise ambiental, propor sujeitos de mudança diferentes e buscar projetos sociais diferentes.

Layrargues e Lima (2014) chamam suas categorias de \emph{macrotendências político-pedagógicas}. Elas são tipos ideais: servem para analisar e comparar, mas não devem funcionar como rótulos rígidos para pessoas, instituições ou projetos. Na prática, pode haver sobreposições. A pergunta adequada não é apenas ``a qual vertente isto pertence?'', mas: \emph{como o problema é explicado, quem é chamado a agir, por quais meios e em favor de qual transformação?}

\section{Questão orientadora e tese geral}

\textbf{Questão orientadora:} quais diagnósticos da crise socioambiental, sujeitos de mudança e projetos de sociedade estão em disputa quando se fala em Educação Ambiental?

A tese que organiza o material é a seguinte: toda prática ambiental educa, mas não educa para o mesmo projeto de sociedade. A perspectiva conservacionista tende a privilegiar natureza, sensibilização e preservação; a pragmática, gestão, consumo e resultados; a crítica, relações de poder, desigualdades, participação e transformação social. A Educação Ambiental Popular aproxima-se do campo crítico, mas guarda ênfases próprias, relacionadas à educação popular, às práticas sociais e à emancipação.

\section{Breve contexto histórico}

O contexto histórico é necessário para entender por que a EA passou a comportar linguagens e projetos diversos. A cronologia abaixo não descreve uma evolução linear, mas a ampliação e institucionalização do debate.

\begin{longtable}{p{0.16\textwidth}p{0.37\textwidth}p{0.37\textwidth}}
\toprule
\textbf{Marco} & \textbf{Impacto para o pensamento ambiental} & \textbf{Relação com a Educação Ambiental} \\
\midrule
\endhead
1972 -- Conferência de Estocolmo & O ambiente ganha centralidade na agenda internacional. O Princípio 19 menciona educação de jovens e adultos e conduta responsável. & Legitima a educação como resposta à questão ambiental, sem fixar uma das vertentes brasileiras posteriores. \\
1977 -- Tbilisi & A formulação de ambiente é ampliada: natural, construído, tecnológico e social; destaca-se a interdependência. & Oferece uma referência para interdisciplinaridade e leitura socioambiental. \\
1981 -- Política Nacional do Meio Ambiente & A Lei n.\textsuperscript{o} 6.938 inclui educação ambiental e participação da comunidade entre princípios da política ambiental brasileira. & Marca a institucionalização nacional inicial. \\
1992 -- Rio-92 e Agenda 21 & Ambiente e desenvolvimento sustentável passam a ser tratados de forma articulada; o capítulo 36 aborda educação, consciência pública e formação. & É contexto decisivo da disputa em torno de ``sustentabilidade''. \\
1999 e 2002 -- PNEA e regulamento & A Lei n.\textsuperscript{o} 9.795 estabelece a EA como componente essencial e permanente, formal e não formal; o Decreto n.\textsuperscript{o} 4.281 regulamenta sua execução. & A política articula pluralismo, participação, cidadania, justiça social e transversalidade, sem adotar uma macrotendência única. \\
\bottomrule
\end{longtable}

Esse percurso ajuda a entender a mudança de foco: de uma preocupação mais concentrada em natureza e conservação, para formulações que incluem sociedade, desenvolvimento, políticas públicas, participação e desigualdades. As vertentes contemporâneas divergem sobre o que essas palavras exigem na prática.

\section{Bases políticas, conceituais e educacionais da EA}

O capítulo de Pelicioni e Philippi Jr. sustenta uma base importante para o restante do debate: EA não é sinônimo de Ecologia e tampouco se reduz à transmissão de informação sobre a natureza. Problemas ambientais possuem dimensões socioeconômicas, políticas e culturais; por isso, sua compreensão requer diálogo entre áreas e um diagnóstico da situação antes de definir objetivos educativos.

Nos artefatos consultados, os seguintes fundamentos foram localizados:

\begin{itemize}
  \item \textbf{Conceitual e epistemológico:} a EA é mais ampla que o conhecimento ecológico; deve considerar as causas sociais, econômicas, políticas e culturais dos problemas.
  \item \textbf{Educacional:} interdisciplinaridade e transdisciplinaridade ajudam a superar a fragmentação do saber. Conhecimento e habilidades devem relacionar-se a atitudes e valores.
  \item \textbf{Ético:} justiça social e responsabilidade diante da vida não são adereços; orientam o sentido da ação educativa.
  \item \textbf{Político:} cidadania inclui participação, representatividade e incidência em políticas públicas; não equivale somente a comportamento privado correto.
  \item \textbf{Práxis:} reflexão que desemboca em ação e retorna à reflexão. Consciência sem ação transformadora é insuficiente.
\end{itemize}

Essa base não apresenta, por si, uma taxonomia fechada de vertentes. Ela torna inteligível por que as classificações posteriores perguntam não apenas ``o que ensinar'', mas também quais relações sociais produzem a crise, quem participa de decisões e qual noção de sustentabilidade está em jogo.

\section{As macrotendências de Layrargues e Lima}

\subsection{O sentido de macrotendência}

Layrargues e Lima interpretam a EA brasileira como um campo social plural e disputado. As macrotendências conservacionista, pragmática e crítica coexistem e competem por hegemonia. Portanto, a classificação é didática, analítica e política; ela não pretende eliminar a complexidade das práticas reais.

\subsection{Conservacionista}

Na macrotendência conservacionista, a questão ambiental aparece com forte ênfase ecológica. Seu horizonte é despertar uma sensibilidade humana para a natureza: conhecer, valorizar, amar e preservar. São recorrentes a biodiversidade, unidades de conservação, biomas, ecoturismo, vivências, trilhas interpretativas e senso-percepção. Santos e Toschi associam essa vertente à ``pauta verde'', à mudança comportamental e a exemplos como trilhas e atividades de contato com a natureza.

O limite apontado por Layrargues e Lima não é que o cuidado ecológico seja irrelevante. A crítica é que ele pode ficar isolado das dimensões sociais, políticas, culturais e de classe que também constituem a crise ambiental. Uma ação pode sensibilizar para a preservação e, ao mesmo tempo, deixar sem resposta quem decide o uso do território, quem lucra e quem sofre os danos.

\subsection{Pragmática}

A macrotendência pragmática enfatiza resultados, gestão e soluções praticáveis. Abrange educação para o desenvolvimento sustentável e para o consumo sustentável, ambientalismo de resultados e ecologismo de mercado. Seus temas frequentes são resíduos, coleta seletiva, consumo consciente, certificação, ecoeficiência, economia verde e responsabilidade socioambiental empresarial.

Santos e Toschi descrevem nela a presença de mídia, empresas e a fórmula ``cada um faz sua parte''. Layrargues e Lima a relacionam ao contexto neoliberal, à redução do Estado e à procura de respostas dentro de limites de mercado. A crítica dos autores é que a ação gerencial, tecnológica ou individual pode ser despolitizada quando não investiga causas e consequências articuladas, relações de poder e as estruturas que produzem danos.

Isso não significa que reduzir desperdício ou organizar a coleta seletiva seja inútil. A questão é se essas ações se encerram em metas e comportamento individual, ou se também abrem discussão sobre produção, consumo, trabalho, desigualdade, políticas públicas e responsabilidade empresarial.

\subsection{Crítica}

A macrotendência crítica reúne correntes nomeadas como Educação Ambiental Popular, emancipatória, transformadora e educação no processo de gestão ambiental. Ela revisa criticamente fundamentos de dominação, mecanismos de acumulação do capital, desigualdade e injustiça socioambiental. Em vez de tratar a crise como simples soma de maus hábitos, procura contextualizar e politizar o debate, problematizando contradições dos modelos de desenvolvimento e de sociedade.

Seus conceitos recorrentes são cidadania, democracia, participação, emancipação, conflito, justiça ambiental, relações de poder e transformação social. A educação é entendida como prática que pode formar sujeitos capazes de interpretar e intervir coletivamente, não apenas como repasse de informação ou treinamento de condutas.

\section{Santos e Toschi: mapeamento didático das vertentes}

Santos e Toschi apresentam a mesma tripartição geral, mas com função mais introdutória e histórico-didática. Seu artigo declara apoio principalmente em Layrargues e Lima (2011), e não deve ser usado para atribuir automaticamente a Layrargues e Lima (2014) exemplos ou formulações específicos.

Sua contribuição ajuda a reconhecer expressões práticas das vertentes. Na conservacionista aparecem trilhas, ecoturismo, agroecologia e sensibilização. Na pragmática, consumo sustentável, resíduos, economia verde, empresas e mídia. Na crítica, educação popular, Paulo Freire, interdisciplinaridade, injustiça ambiental e problematização de causas sociais.

Uma observação relevante das autoras é que uma abordagem inicialmente pragmática pode incorporar dimensão crítica se passar a considerar análises sociais, econômicas, culturais e políticas. Isso reforça que a análise deve observar a coerência entre diagnóstico, objetivos, método e efeitos, não apenas o tema superficial da atividade.

\section{Educação Ambiental Crítica segundo Carvalho}

Carvalho parte da multiplicidade de nomes presentes no campo. Para ela, ``ambiental'' já é uma adjetivação de educação; ``crítica'' é uma segunda marca, que explicita determinada maneira de compreender e fazer educação. A autora não apresenta a EA crítica como a única forma legítima de EA. Sua proposta é tornar legíveis os pressupostos, valores e posições das diferentes formas de nomear a prática educativa.

O conceito de \textbf{endereçamento}, tomado dos estudos de cinema, pergunta como uma prática educativa se constitui e a quem ela se dirige. Não existe educação neutra ou produzida fora de relações sociais: públicos, discursos e instituições são formados em dinâmicas de poder e contrapoder; o destinatário também participa da constituição do artefato educativo.

Carvalho localiza a perspectiva crítica em raízes democráticas e emancipatórias e aproxima-a da tradição brasileira de educação popular. A especificidade ambiental está em compreender relações sociedade--natureza e intervir sobre problemas e conflitos ambientais.

O \textbf{sujeito ecológico} é uma noção central. Trata-se de uma subjetividade solidária com os meios social e ambiental, capaz de identificar, problematizar e agir diante de questões socioambientais, sob um horizonte de justiça ambiental. A autora evita dois reducionismos: não basta somar gestos individuais de pessoas que ``fazem sua parte'', mas também não basta invocar abstratamente o sistema social e adiar a formação da subjetividade. A transformação incide na relação indivíduo--sociedade.

Entre as pretensões da EA crítica aparecem leitura multidimensional dos problemas, mudança de padrões de uso e distribuição de bens ambientais, atitude ética, estética e política, cidadania ambiental, vínculo entre escola e ambientes locais e educador como mediador de relações socioeducativas.

\section{Educação Ambiental Crítica segundo Guimarães}

Guimarães explica por que acrescenta o termo ``crítica'' à EA: para diferenciar uma ação educativa que possa contribuir para transformar uma realidade historicamente marcada por grave crise socioambiental. Para ele, EA crítica é uma contraposição à EA conservadora, e não seu simples aperfeiçoamento.

O autor associa a EA conservadora a cientificismo cartesiano, antropocentrismo, visão fragmentada, individualismo e prática comportamentalista. Ela tende a reproduzir os referenciais que constituem a própria crise, conservando a lógica dominante. A expressão \textbf{socioambiental} é importante justamente porque recusa separar os problemas sociais dos ambientais.

Na perspectiva crítica, a realidade é constituída por forças, conflitos, consensos e relações de poder. Guimarães menciona a Teoria Crítica, Paulo Freire, Milton Santos e Edgar Morin como influências para leitura crítica, espacial e complexa da realidade. Contudo, desvelar conflitos não basta: é preciso \textbf{práxis}, isto é, um movimento em que a reflexão subsidia uma prática criativa e a prática realimenta a reflexão.

A metáfora da ``contra-correnteza'' expressa que a transformação não se obtém pela resistência isolada. Ela demanda uma ação coletiva capaz de enfrentar a corrente dominante. A finalidade educativa é desvelar relações de poder e mecanismos ideológicos, instrumentalizando inserção política e mobilização coletiva para uma sociedade ambientalmente sustentável.

\section{Educação Ambiental Popular segundo Souza}

Souza apresenta a Educação Ambiental Popular como corrente teórica voltada a orientar práticas sociais, escolares e não escolares. Pelo resumo do artigo, confirmado nos artefatos, a EAP pode orientar uma práxis de libertação, emancipação e transformação de realidades injustas, opressoras e excludentes na cotidianidade latino-americana. O artigo explora relações históricas entre EA e educação popular e sua convergência com a macrotendência crítica.

\textbf{Limite de confirmação:} nesta pesquisa, o PDF do artigo não foi extraível de forma confiável. Por isso, não foram atribuídos ao autor exemplos internos, páginas específicas ou uma descrição detalhada de suas relações com Paulo Freire, movimentos sociais e diálogo. Esses vínculos são plausíveis e coerentes com o resumo, mas devem ser confirmados pela leitura integral antes de serem apresentados como formulações textuais de Souza.

Como interpretação integrada, a EAP pode ser apresentada como tradição relacionada ao campo crítico, e não como sinônimo de ``educação para pessoas pobres''. O adjetivo ``popular'' indica uma orientação político-pedagógica ligada a emancipação, participação e práticas sociais.

\section{Comparação geral}

\begin{longtable}{p{0.16\textwidth}p{0.19\textwidth}p{0.19\textwidth}p{0.19\textwidth}p{0.19\textwidth}}
\toprule
\textbf{Critério} & \textbf{Conservacionista} & \textbf{Pragmática} & \textbf{Crítica} & \textbf{Popular} \\
\midrule
\endhead
Natureza & Objeto de cuidado e preservação. & Recurso a ser gerido de modo sustentável. & Relação sociedade--natureza. & Território e vida vividos coletivamente\footnotemark. \\
Crise & Degradação ecológica. & Ineficiência, resíduos e escolhas de consumo. & Desigualdade, desenvolvimento, dominação e poder. & Opressão e injustiça na vida cotidiana\footnotemark[1]. \\
Objetivo & Sensibilizar e preservar. & Melhorar resultados e condutas. & Transformar relações socioambientais. & Emancipar e transformar práticas sociais\footnotemark[1]. \\
Indivíduo & Agente de cuidado. & Consumidor ou gestor responsável. & Sujeito situado, não isolado. & Sujeito dialógico e coletivo\footnotemark[1]. \\
Coletividade/Estado & Papel secundário ou variável. & Parceria, governança e gestão. & Participação, políticas e conflito. & Organização comunitária e participação\footnotemark[1]. \\
Metodologia & Vivência ecológica e sensibilização. & Campanhas, indicadores e gestão. & Problematização, diálogo e práxis. & Diálogo e práticas sociais\footnotemark[1]. \\
Sustentabilidade & Conservação. & Eficiência e mercado. & Justiça socioambiental. & Emancipação e justiça\footnotemark[1]. \\
\bottomrule
\end{longtable}
\footnotetext{Síntese de trabalho. No caso da coluna ``Popular'', os detalhes devem ser confirmados pela leitura integral de Souza (2018).}

\section{Exemplo: resíduos em uma universidade}

O mesmo problema permite visualizar diferenças entre as orientações.

\begin{description}[style=nextline,leftmargin=3.5cm,font=\normalfont\bfseries]
\item[Conservacionista] Pergunta como os resíduos afetam fauna, flora e recursos. Pode usar visita guiada, identificação de materiais, sensibilização e cuidado.
\item[Pragmática] Pergunta como reduzir volume e elevar a reciclagem. Pode realizar auditoria de lixeiras, estabelecer metas, melhorar sinalização, coleta seletiva e consumo responsável.
\item[Crítica] Pergunta quem produz, decide e recebe os danos dos resíduos. Pode mapear contratos, orçamento, cadeia de reciclagem, condições de trabalho dos catadores e desigualdades territoriais; promove debate e incidência por redução na fonte e trabalho digno.
\item[Popular] Pergunta quais saberes e problemas de estudantes, trabalhadores e comunidade devem orientar a ação. Pode construir círculos de diálogo, diagnóstico participativo e ação com os grupos afetados.
\end{description}

Os exemplos não são receitas rígidas. Eles mostram que a diferença está na pergunta de partida e na escala de transformação pretendida. O mesmo raciocínio pode ser aplicado a água, queimadas, consumo, áreas verdes, mudanças climáticas ou conflitos territoriais.

\section{Glossário de estudo}

\begin{description}[style=nextline,leftmargin=4.1cm,font=\normalfont\bfseries]
\item[Educação Ambiental] Processo educativo permanente que articula ambiente, conhecimento, valores e ação; não se reduz à Ecologia.
\item[Macrotendência] Tipo ideal amplo que reúne orientações político-pedagógicas; não é classificação rígida de toda prática.
\item[Pragmatismo ambiental] Ênfase em resultados, gestão, consumo, resíduos, ecoeficiência e soluções compatíveis com o mercado.
\item[Práxis] Relação reflexiva entre ação e reflexão que transforma ambas; não é aplicação mecânica de teoria.
\item[Emancipação] Ampliação de autonomia e capacidade coletiva de intervir em relações opressoras.
\item[Sustentabilidade] Conceito disputado: pode significar adequação gerencial ou justiça e transformação, conforme o referencial.
\item[Cidadania ambiental] Participação em direitos, deveres e decisões públicas que afetam ambiente e coletividade.
\item[Sujeito ecológico] Em Carvalho, subjetividade solidária capaz de agir sobre questões socioambientais.
\item[Justiça ambiental] Atenção à distribuição desigual de bens, danos e poder de decisão ambiental.
\item[Relações de poder] Disputas que constituem sentidos, territórios, decisões e conflitos socioambientais.
\end{description}

\section{Controvérsias e cuidados de leitura}

\begin{enumerate}
\item \textbf{Categorias não são essências.} Um projeto pode combinar dimensões de mais de uma vertente. Deve-se examinar seus pressupostos e efeitos.
\item \textbf{A pragmática não é automaticamente ilegítima.} A crítica dirige-se aos seus limites quando ações de gestão e consumo são despolitizadas. Elas podem ganhar criticidade se articuladas a análise social, econômica, cultural e política.
\item \textbf{EA crítica não é só conscientização.} Em Carvalho e Guimarães, envolve poder, conflito, subjetividade, participação, práxis e ação coletiva.
\item \textbf{EA popular não é uma educação definida pelo público pobre.} Ela remete a uma tradição político-pedagógica de diálogo, emancipação e práticas sociais.
\item \textbf{Sustentabilidade não é neutra.} A questão analítica é: sustentável para quem, por quais meios e com qual distribuição de poder e bens?
\end{enumerate}

\section{Síntese final}

Classificar vertentes não serve para decorar uma lista de definições. Serve para tornar visíveis decisões que costumam ficar implícitas: que natureza está sendo considerada, como a crise é explicada, quem aparece como responsável, qual o lugar do Estado e da coletividade, e que transformação se espera da educação.

A vertente conservacionista ressalta a preservação e a sensibilização; a pragmática, respostas gerenciais e de consumo; a crítica, relações estruturais, poder, conflito, participação e transformação. Carvalho acrescenta a necessidade de enxergar endereçamentos e formar sujeitos ecológicos; Guimarães enfatiza práxis e mobilização coletiva; Souza aproxima o debate da educação popular e da emancipação. A melhor leitura do conjunto preserva essas aproximações sem apagar diferenças de nomenclatura e de ênfase.

\section{Referências}

\begin{thebibliography}{99}
\bibitem{brasil1981} BRASIL. \textit{Lei n.\textsuperscript{o} 6.938, de 31 de agosto de 1981}. Política Nacional do Meio Ambiente. Disponível em: \url{https://www.planalto.gov.br/ccivil_03/leis/l6938compilada.htm}.
\bibitem{brasil1999} BRASIL. \textit{Lei n.\textsuperscript{o} 9.795, de 27 de abril de 1999}. Política Nacional de Educação Ambiental. Disponível em: \url{https://www.planalto.gov.br/ccivil_03/leis/l9795.htm}.
\bibitem{brasil2002} BRASIL. \textit{Decreto n.\textsuperscript{o} 4.281, de 25 de junho de 2002}. Disponível em: \url{https://www.planalto.gov.br/ccivil_03/decreto/2002/d4281.htm}.
\bibitem{carvalho2004} CARVALHO, Isabel Cristina de Moura. Educação ambiental crítica: nomes e endereçamentos da educação. In: LAYRARGUES, Philippe Pomier (coord.). \textit{Identidades da educação ambiental brasileira}. Brasília: MMA, 2004. p. 13--24.
\bibitem{guimaraes2004} GUIMARÃES, Mauro. Educação ambiental crítica. In: LAYRARGUES, Philippe Pomier (coord.). \textit{Identidades da educação ambiental brasileira}. Brasília: MMA, 2004. p. 25--34.
\bibitem{layrargueslima2014} LAYRARGUES, Philippe Pomier; LIMA, Gustavo Ferreira da Costa. As macrotendências político-pedagógicas da educação ambiental brasileira. \textit{Ambiente \& Sociedade}, v. 17, n. 1, p. 23--40, 2014. DOI: \url{https://doi.org/10.1590/1809-44220003500}.
\bibitem{onu1972} ORGANIZAÇÃO DAS NAÇÕES UNIDAS. \textit{Declaração de Estocolmo sobre o Meio Ambiente Humano}. 1972.
\bibitem{onu1992} ORGANIZAÇÃO DAS NAÇÕES UNIDAS. \textit{Agenda 21}, cap. 36: Promoting Education, Public Awareness and Training. 1992. Disponível em: \url{https://www.un.org/esa/dsd/agenda21/res_agenda21_36.shtml}.
\bibitem{pelicioni} PELICIONI, Maria Cecília Focesi; PHILIPPI Jr., Arlindo. Bases políticas, conceituais, filosóficas e ideológicas da Educação Ambiental. In: PHILIPPI Jr.; PELICIONI (org.). \textit{Educação Ambiental e Sustentabilidade}. Barueri: Manole, 2005. p. 3--12. \textit{A edição/ano deve ser conferida no exemplar da disciplina.}
\bibitem{santostoschi2015} SANTOS, Jéssica de Andrade; TOSCHI, Mirza Seabra. Vertentes da Educação Ambiental: da conservacionista à crítica. \textit{Fronteiras: Journal of Social, Technological and Environmental Science}, v. 4, n. 2, p. 241--250, 2015. DOI: \url{https://doi.org/10.21664/2238-8869.2015v4i2.p241-250}.
\bibitem{souza2018} SOUZA, Tiago Zanquêta de. A educação ambiental popular: contribuições em práticas sociais. \textit{Motricidades}, v. 2, n. 1, p. 60--70, 2018. DOI: \url{https://doi.org/10.29181/2594-6463-2018-v2-n1-p60-70}.
\end{thebibliography}

\section{Pendências de verificação}

\begin{itemize}
\item Confirmar no exemplar utilizado pela disciplina se a referência de Pelicioni e Philippi Jr. datada de 2014 é reimpressão ou edição diferente. A pesquisa encontrou com segurança a edição de 2005, p. 3--12.
\item Consultar o PDF integral de Souza (2018) antes de usar páginas internas, exemplos específicos ou citações literais da obra.
\item Conferir diretamente nos PDFs paginados qualquer citação entre aspas. Este resumo usa predominantemente paráfrases rigorosas.
\end{itemize}

\end{document}
